\section{Engineering Reference}
\label{sec:engineering-reference}

This appendix collects implementation-oriented material that supports the main
algorithmic narrative without interrupting its progression.  It covers codebase
organization, template specialization, validation tools, kernel inventory,
display output, PNG encoding, orbit interchange, and cross-platform constraints.

\subsection{Codebase Organization}
\label{sec:code-organization}

The source tree separates operating-system integration, reusable numeric and
rendering code, application front ends, and validation tools.  The build-target
boundaries are significant: they keep CUDA compilation isolated, prevent
platform APIs from leaking into portable interfaces, and allow the Windows GUI,
Linux GUI, and command-line renderer to share the same rendering implementation.

\subsubsection{Dependency layers}

\Cref{fig:codebase-layers} summarizes the architectural layers.  Arrows point
toward targets that consume the layer below them; final executables link the
libraries required by the features they expose.

\begin{figure}[!htbp]
  \centering
  \begin{tikzpicture}[
      component/.style={
        draw,
        rounded corners,
        align=center,
        minimum height=12mm,
        text width=0.90\textwidth,
        inner sep=3mm
      },
      flow/.style={-{Latex[length=2.5mm]}, thick}
    ]
    \node[component] (foundation) {
      \textbf{Foundation libraries}\\
      \code{FractalSharkPlatform}\qquad \code{HpSharkFloatLib}
    };
    \node[component, above=8mm of foundation] (rendering) {
      \textbf{Shared rendering implementation}\\
      \code{FractalSharkLib}\qquad \code{FractalSharkGpuLib}
    };
    \node[component, above=8mm of rendering] (consumers) {
      \textbf{Applications and developer tools}\\[1mm]
      \code{FractalShark} + \code{FractalSharkGuiWin32}\quad
      \code{FractalSharkGuiLinux}\quad \code{FractalSharkCli}\\
      \code{FractalTray}\quad \code{FractalSaver}\quad
      \code{FractalSharkTest}\quad \code{HpSharkFloatTest}\\
      \code{HpSharkFloatTestLib}\quad \code{HpSharkInstantiate}
    };

    \draw[flow] (foundation) -- (rendering);
    \draw[flow] (rendering) -- (consumers);
  \end{tikzpicture}
  \caption[Codebase dependency layers]{Codebase dependency layers.  Platform and
    high-precision arithmetic support the shared CPU and CUDA rendering code;
    applications and validation tools consume those shared layers.}
  \label{fig:codebase-layers}
\end{figure}

\subsubsection{Foundation and rendering libraries}

\paragraph{\code{FractalSharkPlatform}.}
The platform library owns operating-system services used by otherwise portable
code.  Its common interfaces cover allocation, files, process information,
timing, diagnostics, crash handling, wait-cursor state, process memory limits,
and native OpenGL contexts.  Separate Win32 and Linux translation units provide
the implementations described in \cref{subsec:platform-abstraction}.

\paragraph{\code{HpSharkFloatLib}.}
The high-precision GPU arithmetic library implements the limb-based
\code{HpSharkFloat} representation, addition, NTT multiplication,
normalization, periodicity support, and the GPU reference-orbit machinery.  It
also contains the parameter families used to specialize those kernels.  The
numeric representation and execution pipeline are documented in
\cref{subsec:ref-orbit-gpu,subsec:hpfloat-model,sec:hp-add,sec:ntt-multiply}.

\paragraph{\code{FractalSharkLib}.}
The core library owns portable fractal state and CPU-side orchestration.  Its
responsibilities include coordinate conversion, reference-orbit dispatch,
linear-approximation construction, feature finding, palettes, iteration-buffer
management, render-thread coordination, persistence, PNG output, shared
commands, and functional tests.  The principal flows appear in
\cref{sec:ref-orbit-calc,sec:render-thread-pool,sec:disk_backed_growable_vectors}.

\paragraph{\code{FractalSharkGpuLib}.}
The rendering CUDA library owns the per-pixel launch surface and kernels for
direct escape-time rendering, perturbation, BLA and LA~v2 traversal,
antialiasing, coloring, and reductions.  It consumes numeric and state types
defined by the core and high-precision libraries; final applications link all
required libraries together.  Kernel behavior is covered in
\cref{sec:mandel-base-kernels,sec:perturb-only,sec:la-v2-perturb,sec:kernel-list}.

\subsubsection{Applications and front ends}

\paragraph{\code{FractalShark} and \code{FractalSharkGuiWin32}.}
The primary Windows application is split across two Visual~Studio projects.
\code{FractalShark} supplies the executable target, resources, icons, and final
link, while \code{FractalSharkGuiWin32} implements the Win32 entry point, main
window, command dispatcher, popup menus, console, saved-location UI, and splash
window.  Shared command identifiers remain in the core library.  User-visible
behavior and frame presentation are described in
\cref{sec:ui-popup-menu,sec:keyboard-shortcuts,subsec:opengl-display}.

\paragraph{\code{FractalSharkGuiLinux}.}
The native Linux application combines an Xlib event loop and menus, GLX context
management, Dear~ImGui overlays, clipboard integration, and the shared renderer.
Its platform boundary and current parity limitations are documented in
\cref{subsec:linux-gui}.

\paragraph{\code{FractalSharkCli}.}
The portable command-line executable provides headless rendering from built-in
views, saved locations, or explicit coordinates. It supports single-shot
renders and a persistent server that accepts requests from other invocations
of the executable. The server reuses a \code{Fractal} instance and calls the
same core and CUDA rendering libraries without a window system. Its command
contract and request flow are documented in \cref{sec:fractalshark-cli}.

\paragraph{\code{FractalTray}.}
The Windows tray application batch-renders a sequence of coordinates through
the shared renderer.  Its input format, configuration, and operating procedure
are described in \cref{subsec:fractaltray}.

\paragraph{\code{FractalSaver}.}
The Windows screen-saver project is legacy, unsupported code retained in the
Visual~Studio solution.  It is not part of the CMake build or release packages
and should not be treated as an active front end.

\subsubsection{Validation and generation tools}

\paragraph{\code{FractalSharkTest}.}
The portable unit-test executable covers core numeric types, coordinate
conversion, linear-approximation data, allocators, serialization, command
state, render goldens, and related utilities.  The custom test framework and
representative coverage are summarized in \cref{subsec:fractalsharktest}.

\paragraph{\code{HpSharkFloatTestLib} and \code{HpSharkFloatTest}.}
The support library contains reference implementations, correctness tests,
performance drivers, and result tracking for high-precision CUDA arithmetic.
The executable supplies the CUDA-aware harness that invokes those tests.
Runtime validation requires physical CUDA hardware, as discussed in
\cref{subsec:hpsharkfloattest}.

\paragraph{\code{HpSharkInstantiate}.}
The standalone developer tool generates explicit-instantiation CUDA sources
for the fixed \code{SharkFloatParams} families.  Generated translation units
control template compile time and final binary coverage; the underlying
strategy is documented in \cref{subsec:template-explicit-inst}.

\subsubsection{Build-system coverage}

The Visual~Studio solution contains the Windows application, Windows GUI
library, tray application, legacy screen saver, shared libraries, tests, and
generation tool.  CMake builds the portable shared libraries and tools,
\code{FractalSharkCli}, and \code{FractalSharkGuiLinux}; it intentionally omits
the Win32-only front ends.  Parallel project-file maintenance, Linux target
composition, and continuous-integration coverage are described in
\cref{subsec:dual-build}.
%% ===================================================================
\subsection{Template design}%
\label{sec:template-design}
%% ===================================================================

\FractalShark{} makes pervasive use of C++ templates to select numeric
types, iteration widths, and algorithmic variants at compile time rather
than at run time.  Two considerations motivate this choice.  First,
the inner loops of both reference-orbit computation and per-pixel
perturbation execute billions of arithmetic operations; virtual dispatch
or run-time branch mispredictions on every iteration would impose
measurable overhead.  Second, encoding precision contracts in the type
system makes numeric conversions and storage formats explicit at template
boundaries, which helps review the deliberate projection from high-precision
recurrence state to reduced perturbation data.

The template parameters fall into four largely orthogonal dimensions,
described in the following subsections.


\subsubsection{Iteration-counter width (\code{IterType})}
\label{subsec:template-itertype}

Every orbit and perturbation result is indexed by an iteration counter
whose type is either \code{uint32\_t} or \code{uint64\_t}.  A 32-bit
counter suffices for the vast majority of renders, supporting up to
$2^{32}-1 \approx 4.3\times10^{9}$ iterations.  Certain ultra-deep
locations, however, require iteration limits that exceed this range.
Switching to a 64-bit counter doubles the per-element storage cost for
every orbit and iteration buffer, so \FractalShark{} treats the counter
width as a compile-time choice rather than unconditionally paying the
wider cost.  The \code{IterType} parameter threads through
\code{PerturbationResults}, \code{LAReference}, and the rendering
kernels, ensuring that all components agree on the iteration index
representation.


\subsubsection{Reference-orbit type (\code{T})}
\label{subsec:template-t}

The parameter~\code{T} selects the numeric type used to store the emitted
reference orbit $\{z_n^\star\}$ and the representation consumed by downstream
perturbation code.  It does not select the recurrence arithmetic: CPU reference
backends use MPIR internally, while the CUDA reference engine uses a
precision-specialized \code{HpSharkFloat} internally and convert emitted samples
to~\code{T}.
The principal choices are:

\begin{itemize}
  \item \code{float} and \code{double}---IEEE~754 single and double
        precision, used at shallow zoom depths where the full mantissa
        suffices to represent the orbit accurately.
  \item \code{CudaDblflt<MattDblflt>}---a double-single representation
        that pairs two 32-bit floats to obtain roughly double the
        mantissa bits of \code{float} alone, with natively vectorisable
        GPU arithmetic (\cref{sec:extended-precision-types}).
  \item \code{HDRFloat<double>} and \code{HDRFloat<float>}---a
        high-dynamic-range wrapper that adds an explicit exponent to the
        mantissa, preventing underflow in the deep-zoom regime where
        orbit values span hundreds of orders of magnitude
        (\cref{sec:hdr-float}).
  \item \code{HDRFloat<CudaDblflt<MattDblflt>>}---the composition of
        the HDR exponent with the double-single mantissa, providing
        both extended range and extended precision.
\end{itemize}

\noindent
The type~\code{T} appears as the second template parameter of
\code{PerturbationResults<IterType,\allowbreak T,\allowbreak PExtras>}
and as the orbit-storage type inside \code{RefOrbitCalc}.


\subsubsection{Perturbation subtype (\code{SubType})}
\label{subsec:template-subtype}

While~\code{T} governs the stored precision of the shared reference orbit,
per-pixel perturbation operates at a lower precision because each
pixel evolves only a small delta $\Delta z_n$ relative to $z_n^\star$.
The \code{SubType} parameter selects this lower-precision type.  It
is typically \code{float}, \code{double}, or
\code{CudaDblflt<MattDblflt>}, and for HDR-wrapped references the
corresponding \code{HDRFloat} variant.  The mapping from \code{T} to
\code{SubType} is encoded by the \code{SubTypeChooser} trait: when
\code{T} is a plain scalar the sub-type is the same scalar; when
\code{T} is an \code{HDRFloat<S>} the sub-type extracts the
underlying~\code{S}.

\code{SubType} appears as an explicit template parameter on
\code{LAReference<IterType,\allowbreak T,\allowbreak SubType,\allowbreak
PExtras>} and is threaded into the perturbation kernels, where it
controls the precision of the per-pixel inner loop.


\subsubsection{Perturbation extras (\code{PExtras})}
\label{subsec:template-pextras}

The \code{Perturb\-Extras} enumeration is a compile-time flag attached to
every \code{Perturbation\-Results} instantiation.  It controls whether
the orbit stores auxiliary data beyond the orbit values themselves:

\begin{itemize}
  \item \code{Disable}---store only the orbit; no extra bookkeeping.
  \item \code{Bad}---record \emph{glitch} information (iterations where
        the perturbation approximation loses accuracy) alongside the
        orbit, enabling downstream glitch-correction passes.
  \item \code{SimpleCompression} and \code{MaxCompression}---store the
        metadata required for reference-orbit compression
        (\cref{sec:ref-compression-zhuoran}).
\end{itemize}

\noindent
Making this a template parameter rather than a run-time flag eliminates
the storage overhead of the auxiliary arrays in the common case where
they are not needed, and allows the compiler to elide the associated
bookkeeping code entirely.


\subsubsection{Instantiation combinatorics}
\label{subsec:template-combinatorics}

The Cartesian product of the four dimensions yields a large number of
potential type combinations.  With two \code{IterType} widths, six
representative \code{T} values, and three commonly used
\code{PExtras} settings (\code{Disable}, \code{Bad},
\code{Simple\-Compression}), the \code{PerturbationResults} class
alone accounts for $2\times 6\times 3 = 36$ concrete instantiations.
This count is visible in the \code{AwesomeVariant} type alias defined
in \code{RefOrbitCalc.h}, which enumerates all 36~alternatives as a
\code{std::variant} so that perturbation results of any type
combination can be stored and dispatched uniformly at run time.

Additional template parameters (\code{SubType}, benchmark modes, reuse
modes) on related classes further multiply the instantiation count.  The
full set of explicit instantiations is available in
\code{RefOrbitCalcTemplates.h}.


\subsubsection{Explicit instantiation strategy}
\label{subsec:template-explicit-inst}

Because every translation unit that includes a template header would
otherwise independently instantiate each specialization it references,
\FractalShark{} relies on \emph{explicit template instantiation} to
centralize code generation and control compile times.

For the perturbation pipeline, the header
\code{RefOrbitCalcTemplates.h} defines a family of preprocessor macros
(\code{InstantiateIsPerturbation\-Result\-UsefulHere},
\code{InstantiateAddPerturbation\-Reference\-Point}, and others) that
expand to explicit instantiation declarations for every supported
\code{(IterType, T, SubType, PExtras)} combination.  Each macro accepts
the type arguments, expands to the required \code{template} declaration,
and is invoked once per combination immediately below its definition.
This pattern confines the instantiation cost to a single translation
unit while allowing other files to use only the declarations.

For the GPU arithmetic subsystem, a separate code-generation tool
(\code{HpShark\-Instantiate}) produces \code{.cu} files containing
explicit instantiations of \code{HpSharkFloat}, the NTT multiply
kernels, the reference-orbit GPU loop, and associated helper templates.
The generator partitions the output into paired
files---one for production parameter sets
(\code{SharkParams1}\,--\,\code{12}),
one for non-periodicity variants
(\code{SharkParamsNP1}\,--\,\code{NP12})---and groups
instantiations into named \emph{batches} (e.g.\
\code{HpSharkFloat\_Conversions}, \code{HpSharkReference},
\code{ReferenceGpuLoop}, and \code{SharkNTT\_Primitives}) so that each batch compiles as an
independent translation unit.  This decomposition enables parallel
compilation and limits the peak memory consumption of the CUDA
compiler.


\subsubsection{The \code{SharkFloatParams} pattern}
\label{subsec:template-sharkparams}

The GPU high-precision arithmetic library takes the compile-time
parameterization one step further by encoding not only type choices but
also \emph{algorithmic configuration} in a template parameter.
The struct
\code{GenericSharkFloatParams<pNumDigits,\allowbreak Periodicity,\allowbreak
NewtonRaphson,\allowbreak SubTypeT>}
bundles several \code{static constexpr} members that the CUDA compiler
can fold into kernel code:

\begin{itemize}
  \item \code{GlobalNumUint32}---the number of 32-bit limbs in each
        high-precision operand.  Despite its name, \code{pNumDigits} is this
        limb count, not a bit or decimal-digit count.  Because
        the compiler knows this value at template instantiation time,
        loop trip counts in the NTT multiply and carry-propagation
        kernels are compile-time constants, enabling aggressive
        unrolling and register allocation.
  \item \code{EnablePeriodicity}---controls whether the reference-orbit
        kernel checks for period detection at each iteration.  Disabling
        it in the non-periodicity (\code{SharkParamsNP*}) and
        Newton--Raphson (\code{SharkParamsNR*}) variants eliminates
        the associated branch and storage overhead.
  \item \code{EnableNewtonRaphson}---when set, the kernel additionally
        tracks the derivative $d_n = \partial z_n/\partial c$ alongside
        the orbit, allocating extra GPU scratch memory and executing the
        derivative recurrence at each iteration.  This derivative is required by
        the Feature Finder's Newton's method solver
        (\cref{sec:feature-finder}).
  \item \code{SharkNTT::Plan}---the compile-time base plan constructed from
        \code{GlobalNumUint32} by \code{SharkNTT::BuildPlan}.  Setup may allocate
        and cache a larger active plan so exponent-shifted polynomial support
        can select a transform length at run time.
\end{itemize}

\noindent
Concrete aliases such as
\code{SharkParams1}\,--\,\code{SharkParams12} fix
\code{pNumDigits} to powers of two from $256$ through $524{,}288$
32-bit limbs.  The production range is therefore 8,192 through 16,777,216
bits, or approximately 2,466 through 5,050,445 decimal digits, with
periodicity enabled and Newton--Raphson disabled.
Additional alias families---\code{SharkParamsNP*} (no periodicity),
\code{SharkParamsNR*} (Newton--Raphson enabled),
\code{SharkParamsDbl*} (\code{double} sub-type), and
\code{SharkParamsDbf*} (\code{CudaDblflt} sub-type)---provide the
remaining feature-flag and sub-type permutations.  Non-production and test
families also include smaller 64- and 128-limb cases, so their alias sets should
not be inferred from the twelve production buckets.  The code generator in
\cref{subsec:template-explicit-inst} emits only the combinations requested by
each kernel batch.

The remaining sections collect operational engineering references for tests,
kernels, display output, orbit interchange, and portability.
% -----------------------------------------------------------------
\subsection{Functional Test Suite}
\label{subsec:crummy-test}

\code{CrummyTest} is a self-contained functional test suite invoked
from the application's right-click context menu
(\code{IDM\_BASICTEST}, labeled ``Run Basic Test'').  The suite
exercises the direct rendering path
(\cref{subsec:two-render-paths}): each test calls \code{Drain()}
first to ensure no pool workers remain active, then invokes
\code{CalcFractal(true)}, which writes results to \code{m\_CurIters}.

The top-level \code{TestAll()} method runs the full suite in sequence:
\begin{itemize}
  \item \code{TestWindowResize}---resizes the application window and
        verifies that a re-render completes correctly.
  \item \code{TestBasic}---iterates through every
        \code{RenderAlgorithm} entry, rendering the default view,
        View~5, and View~11 with both 32-bit and 64-bit iteration
        types.
  \item \code{TestReferenceSave}---saves and reloads reference orbits,
        verifying round-trip correctness.
  \item \code{TestVariedCompression}---tests reference-orbit
        compression at varying error thresholds.
  \item \code{TestImaginaLoad}---loads orbits stored in the Imagina
        format (\cref{subsec:imagina-format}).
  \item \code{TestStringConversion}---validates string-to-precision
        and precision-to-string conversion utilities.
  \item \code{TestPerturbedPerturb}---exercises perturbed iteration
        chains for numerical correctness.
  \item \code{TestGrowableVector}---validates the
        \code{GrowableVector} data structure
        (\cref{sec:disk_backed_growable_vectors}).
\end{itemize}
A separate \code{TestReallyHardView27} (not included in
\code{TestAll}) renders View~\#27, a period ${\sim}28$~billion point
requiring hours of computation; the test serves as an overnight
stress benchmark rather than a routine regression check.

% -----------------------------------------------------------------
\subsection{PNG Encoder and Save-Path Validation}
\label{subsec:png-validation}

The PNG encoder described in \cref{subsec:png-output} has a focused
hardware suite in \code{FractalSharkCudaTest}. The \code{CudaPng\_*}
cases encode identical input through the GPU encoder and the CPU's
current PNG save settings. LodePNG decodes
both outputs to RGBA16 for exact comparison with independently
serialized expected pixels. Coverage includes low channel bits,
transparency, padded strides, row and region boundaries, deterministic
output, noisy input, multiple IDAT chunks, workspace reuse, concurrent
streams, and invalid inputs.

Internal filtering tests reconstruct scanlines independently. Raw zlib
tests cover length/distance coding, overlapping copies, stored blocks,
checksums, and final-block termination. Huffman tests compare code
lengths and weighted cost with the bundled CPU implementation, including
degenerate alphabets, stable ties, and skewed frequencies requiring
length limiting. Header tests independently expand repetition codes.
Corruption tests separate PNG CRC failures from Adler-32 failures by
repairing the enclosing chunk CRC. Internal test interfaces in
\code{PngCompressionKernels.cuh} use production device routines and
also expose the fixed/stored baseline for size comparisons.

The \code{CudaPngSave\_*} cases exercise the production dispatch path:
both iteration widths, AA factors~1--4, odd dimensions, repeated saves
and resizing, copied/moved snapshots, render-pool frame publication,
synchronous Background completion, Unicode filenames, existing-file
skips, write failures, and switching back to queued CPU saves. Rotated
Default and custom palettes are compared exactly with the CPU saver for
both iteration widths and AA factors~1--4. Save snapshots are constructed
before subsequent rotation and auxiliary-depth mutations to verify that
captured settings survive renderer reuse.

The \code{CudaRecolor\_*} cases compare host-buffer GPU recoloring and
normal GPU frame extraction against independent expected colors. Coverage
includes padded input rows, odd dimensions, both iteration widths, AA
factors~1--4, interior samples, rotation saturation and overflow boundaries,
Default and custom palette tables, Basic grayscale, scalar-setting reuse,
output-capacity checks, and recovery after invalid requests. Queued tests
exercise repeated rotation, reset, resizing, paced cancellation, disabled
repainting, later rendering, and unchanged iteration counts.
\code{FractalSharkLib\_RecolorCpuWithoutCuda} exercises the same queued
lifecycle with CUDA disabled. The existing CPU save-pool and PNG tests
also remain runnable without CUDA hardware.

Run the following PowerShell commands from the repository root after
building both Windows configurations:

\noindent\begin{minipage}{\linewidth}
\begin{lstlisting}[language={},upquote=true]
foreach ($configuration in 'Debug', 'Release') {
    & ".\$configuration\FractalSharkCudaTest.exe" --use-gpu `
        --filter=CudaPng_* --filter=CudaPngSave_* --filter=CudaRecolor_*
    & ".\$configuration\FractalSharkTest.exe" `
        --filter=FractalSharkLib_SavePool* `
        --filter=FractalSharkLib_PngEncoding* `
        --filter=FractalSharkLib_RecolorCpuWithoutCuda
}
\end{lstlisting}
\end{minipage}

\paragraph{Existing render goldens.}
Two existing View0 cases exercise CPU64 and GPU1x32 PNG saves through
\code{SaveCurrentFractal}. Their fixtures drain the render pool, call
\code{CalcFractal(true)}, save the current frame, and wait for completion.
The eight existing GPU color goldens additionally cover both iteration widths
and AA factors~1--4. Only these ten cases are selected by the following commands:

\noindent\begin{minipage}{\linewidth}
\begin{lstlisting}[language={},upquote=true]
foreach ($configuration in 'Debug', 'Release') {
    & ".\$configuration\FractalSharkTest.exe" --use-gpu `
        --filter "RenderGolden_Basic_Cpu64_View0_Bits32_Store0_*" `
        --filter "RenderGolden_Basic_Gpu1x32_View0_Bits32_Store0_*" `
        --filter "RenderGolden_GpuColors_*"
}
\end{lstlisting}
\end{minipage}

Render goldens compare decoded RGBA16 pixels, so different compression
bytes alone require no golden update. Intentional pixel changes use the
selected-case update procedure in \cref{subsec:render-regression-goldens};
unaffected cases need not be regenerated or rerun.

\paragraph{Benchmark procedure.}
The standalone benchmark uses opaque uniform, gradient, seeded-noise,
and GPU-rendered Mandelbrot images at 3840~by~2160 resolution:

\begin{lstlisting}[language={},upquote=true]
.\Release\FractalSharkCudaTest.exe --use-gpu --filter=CudaPngBenchmark_4K
\end{lstlisting}

Run it after other builds and tests complete. Every output is decoded
and checked exactly. Compression size must not exceed the fixed/stored
baseline, and uniform, gradient, and Mandelbrot fixtures must improve
on that baseline. Timings are diagnostic rather than assertions;
\cref{subsec:png-benchmarks} defines the measurement boundaries and
reports the recorded results.

\paragraph{Compute Sanitizer on Windows.}
Select named-pipe transport before launching memcheck:

\begin{lstlisting}[language={},upquote=true]
$env:NV_COMPUTE_SANITIZER_LOCAL_CONNECTION_OVERRIDE = 'named-pipes'
& "$env:CUDA_PATH/compute-sanitizer/compute-sanitizer.exe" `
    --tool memcheck --error-exitcode 99 `
    .\Debug\FractalSharkCudaTest.exe --use-gpu `
    --filter=CudaPng_* --filter=CudaPngSave_* --filter=CudaRecolor_*
\end{lstlisting}

The sanitizer communicates with its instrumented process. Named pipes
keep that Windows communication local without a TCP listener; TCP
transport can otherwise produce a firewall prompt attributed to the
test executable. Ordinary encoder tests start no network service.
The transport override is documented in NVIDIA's
\href{https://docs.nvidia.com/compute-sanitizer/ComputeSanitizer/index.html\#environment-variables}
{Compute Sanitizer environment-variable reference}.

% -----------------------------------------------------------------
\subsection{Selected Kernel Entry Points}
\label{sec:kernel-list}

This section lists representative CUDA entry points.  It is intentionally not
exhaustive: templated helpers, wrapper launch functions, generated
specializations, and internal kernels change more frequently than this appendix.

\begin{itemize}

\item \textbf{Mandelbrot base kernels} (\cref{sec:mandel-base-kernels})
  \begin{itemize}
    \item \code{mandel\_1x\_float}
    \item \code{mandel\_1x\_double}
    \item \code{mandel\_2x\_float}
    \item \code{mandel\_2x\_double}
    \item \code{mandel\_4x\_float}
    \item \code{mandel\_4x\_double}
    \item \code{mandel\_hdr\_float}
  \end{itemize}

\item \textbf{Mandelbrot perturbation kernels}
  \begin{itemize}
    \item \code{mandel\_1x\_double\_perturb\_bla}
    \item \code{mandel\_1xHDR\_float\_perturb\_bla} (\cref{sec:bilinear-approx})
    \item \code{mandel\_1xHDR\_float\_perturb\_lav2} (\cref{sec:la-v2-perturb,sec:perturb-only})
  \end{itemize}

\item \textbf{High-precision GPU reference kernels}
  (\cref{subsec:ref-orbit-gpu,subsec:fused-ntt-reference-step})
  \begin{itemize}
    \item \code{HpSharkReferenceSetupKernel}
    \item \code{HpSharkReferenceGpuLoop}
  \end{itemize}

\item \textbf{Utility kernels}
  \begin{itemize}
    \item \code{antialiasing\_kernel}
    \item \code{max\_reduce}
    \item \code{max\_kernel}
  \end{itemize}

\item \textbf{PNG encoding kernels} (\cref{subsec:png-output})
  \begin{itemize}
    \item Serialization and scanline filter selection
    \item Regional LZ77 encoding and dynamic Huffman construction
    \item Compaction, Adler-32, and PNG chunk assembly/CRC
  \end{itemize}

\item \textbf{Reference NTT primitives} (\cref{sec:ntt-multiply})
  \begin{itemize}
    \item \code{SharkNTT\_Primitives}
  \end{itemize}

\item \textbf{Disabled kernels}
  \begin{itemize}
    \item \code{mandel\_2x\_float\_perturb\_setup}
    \item \code{mandel\_2x\_float\_perturb}
  \end{itemize}

\end{itemize}

\subsubsection{Commentary}

\begin{itemize}

\item All the per-pixel CUDA kernels are in \code{gpu\_render.cu}. The two
most interesting are probably \code{mandel\_\allowbreak 1xHDR\_\allowbreak float\_\allowbreak perturb\_bla} and
\code{mandel\_\allowbreak 1xHDR\_\allowbreak float\_\allowbreak perturb\_lav2}, which are the ones I've spent the
most time on lately. For better performance at low zoom levels, consider
\code{mandel\_1x\_float\_perturb}, which leaves out linear approximation
and just does straight perturbation up to $\sim 10^{30}$, which corresponds
with the 32-bit float exponent range.

\item The \code{mandel\_1x\_float} (\cref{sec:mandel-1x-float}) is the classic 32-bit float Mandelbrot and
is screaming fast on a GPU. This one is optimized with fused multiply-add for
fun even though it is of limited practical use because one can barely zoom in before
pixellation sets in.

\end{itemize}
\subsection{OpenGL Display Pipeline}
\label{subsec:opengl-display}

The GL consumer thread inside \code{RenderThreadPool} dequeues completed
frames from a \code{FrameCompletionQueue} and uploads each frame as an
OpenGL texture.  Before any GL call is issued, the thread makes the
shared context current on itself.

Two rendering paths are supported.  The \emph{software path} targets
GL~1.1--compatible drivers (including the GDI Generic fallback) by tiling
the frame into power-of-two chunks that respect the
\code{GL\_MAX\_TEXTURE\_SIZE} limit---typically 1024 on GDI Generic---and
converting \code{Color16} pixels to RGBA8 before upload.  The
\emph{hardware path} uploads the entire frame as a single
non-power-of-two texture in RGBA16 format via \code{glTexImage2D},
avoiding the tiling overhead entirely.

Both paths flip the Y~axis during texture-coordinate generation to
account for the mismatch between the top-left origin used by the
frame buffer and the bottom-left origin expected by OpenGL.

% -----------------------------------------------------------------
\subsection{PNG Output and GPU Compression}
\label{subsec:png-output}

\FractalShark{} writes lossless 16-bit PNG images through the backend
selected by the resolved render algorithm. CPU algorithms retain
\code{WPngImage} encoding and the background save thread pool. GPU
algorithms use the CUDA encoder and finish encoding and writing on the
thread that initiated the save. AUTO follows the resolved algorithm's
\code{UseLocalColor} setting. Both paths share filename resolution,
extension handling, and the behavior of skipping an existing output file.

\subsubsection{Save dispatch and snapshot ownership}

\code{SaveCurrentFractal} captures the encoder backend, dimensions,
antialiasing factor, iteration type and limit, palette, and iteration
buffer in \code{PngParallelSave}. A copied snapshot owns an independent
iteration buffer. A moved snapshot installs a replacement current buffer
and returns its original storage to the spare-buffer pool after saving.
Queued CPU saves consequently remain valid across later renders or
resizes. Iteration-text saves also retain the background worker path.

Callers drain the render pool before reading the authoritative
\code{m\_CurIters} snapshot (\cref{subsec:two-render-paths}). The GPU save
path initializes renderer~0 from that snapshot and uploads its logical
iteration rows using the padded host row stride. The upload avoids any
dependency on which renderer produced the published frame. Existing
GPU coloring and antialiasing kernels then produce the final packed,
row-major \code{Color16} image on the compute stream. The device image
passes directly to a lazily allocated, reusable PNG encoder; no
uncompressed color image is downloaded for encoding.

GPU saves apply the same palette rotation, saturation, auxiliary-depth shift,
and palette lookup as CPU saves, using the snapshot's captured settings
(\cref{subsec:coloring}).
Encoding is lossless with respect to those final colors. GPU encoding
and file-write failures return directly to the caller, whereas CPU
background failures use the existing worker diagnostics. The GPU bridge
waits for pending stream work before releasing the save snapshot and
does not retain its palette address as a host cache key.

\code{PngCompletionMode::Background} permits background completion only
for CPU saves. GPU saves are complete when the save call returns.
The corresponding CLI response and timing behavior is documented in
\cref{sec:fractalshark-cli}.

\subsubsection{Encoder interface and lifetime}

The CUDA-facing interface is declared in
\code{FractalSharkGpuLib/PngEncoder.cuh}:

\begin{lstlisting}
Png::GpuPngEncoder encoder;
std::vector<unsigned char> pngBytes;
const cudaError_t error = encoder.Encode(
    devicePixels, width, height, rowStrideBytes, stream, pngBytes);
\end{lstlisting}

The input is a caller-owned device allocation of native-endian 16-bit
RGBA samples with an explicit byte stride. Row padding is ignored;
the allocation must cover the requested image and remain valid until
the call returns. Alpha is straight, matching \code{WPngImage}.
Opaque images use RGB16, PNG color type~2; any non-opaque sample selects
RGBA16, color type~6. Low channel bits and transparent pixels' RGB
values are preserved. Output is non-interlaced and contains the PNG
signature followed by IHDR, IDAT, and IEND chunks
\cite{PNGThirdEdition2025}.

\code{Encode} returns \code{cudaSuccess} with completed host bytes,
or a CUDA error with an empty output vector. Host allocation exceptions
propagate. Dimension, stride, alignment, and size-overflow checks precede
kernel execution; accessible device storage remains a caller requirement.
Each encoder context binds to the device used for its first valid call.
Calls on one context must be serialized, while separate contexts and
streams may execute concurrently.

\code{GetWorkspaceBytes()} reports retained device allocation capacity,
excluding the input buffer and host output. Workspace grows as needed
and is reused for smaller images. Destruction releases the
workspace on its owning device. Encoding waits for its supplied stream
to obtain alpha information, compressed sizes, and completed bytes;
it does not call \code{cudaDevice\allowbreak Synchronize}. Workspace allocation
uses \code{cudaMalloc}/\code{cudaFree}, which can impose additional
allocation synchronization.

\subsubsection{Serialization and scanline filtering}

\Cref{fig:png-encoder-pipeline} summarizes the GPU save path.
\code{PngEncoder.cu} serializes samples in big-endian byte order.
Each scanline evaluates the PNG None, Sub, Up, Average, and Paeth
filters using six bytes per RGB16 pixel or eight per RGBA16 pixel
\cite{PNGThirdEdition2025}. The minimum-sum scoring matches the bundled
LodePNG heuristic, including its unsigned None score and
\code{255 - value} convention for negative differences. Equal scores
select the lowest filter number. Only the chosen filtered row is
materialized in the compression input.

\begin{figure}[!htbp]
\centering
\begin{tikzpicture}[
  stage/.style={draw, rounded corners, align=center, font=\small,
                text width=3.7cm, minimum height=1.1cm},
  flow/.style={-{Latex[length=2mm]}, thick},
  node distance=0.55cm and 0.6cm]
  \node[stage] (color) {Upload snapshot\\GPU coloring and AA};
  \node[stage, right=of color] (filter) {Serialize samples\\Select row filters};
  \node[stage, right=of filter] (deflate) {Regional DEFLATE\\encoding};
  \node[stage, below=of deflate] (frame) {GPU checksums\\and PNG assembly};
  \node[stage, left=of frame] (transfer) {Transfer PNG bytes\\to host};
  \node[stage, left=of transfer] (write) {Write output file\\Return save status};
  \draw[flow] (color) -- (filter);
  \draw[flow] (filter) -- (deflate);
  \draw[flow] (deflate) -- (frame);
  \draw[flow] (frame) -- (transfer);
  \draw[flow] (transfer) -- (write);
\end{tikzpicture}
\caption[GPU PNG save pipeline]{GPU PNG save pipeline. Coloring,
filtering, compression, checksums, and PNG assembly use device data.
Only small result metadata and completed PNG bytes return to the host;
the initiating thread waits for completion and writes the file.}
\label{fig:png-encoder-pipeline}
\end{figure}

\subsubsection{Regional DEFLATE and dynamic Huffman coding}

Filtered bytes are divided into independent regions of at most
32~KiB. A warp processes each region with a 4,096-entry hash table,
a greedy LZ77 parser, and warp-assisted match comparisons. Matches
reference only earlier positions within their own region, may overlap,
and have a maximum length of 258~bytes. Ordered position values and
\code{atomicMax} make hash insertion deterministic when positions
consumed by one match collide. Hash tables and intermediate buffers
reside in allocated device workspace.

Each region first emits a non-final fixed-Huffman block followed by an
empty stored block that restores byte alignment. If that representation
is at least as large as a stored block, a non-final stored block replaces
it. A second kernel constructs dynamic Huffman tables and replaces the
candidate only when the complete dynamic representation is strictly
smaller. Equal sizes retain the existing candidate. Block syntax and
length/distance coding follow DEFLATE \cite{RFC1951}.

\code{PngHuffman.cuh} implements iterative boundary package merge for
length-limited Huffman trees: 15~bits for literal/length and distance
codes, and 7~bits for the header alphabet. Stable frequency/symbol
ordering resolves ties deterministically. Empty and single-symbol
alphabets use two one-bit codes for decoder compatibility; the
end-of-block symbol always receives a frequency of one. Canonical
codes are reversed for DEFLATE bit emission.

The dynamic header trims trailing unused symbols and run-length encodes
the combined code lengths with symbols~16, 17, and~18. Candidate size
includes all header and payload bits, the end-of-block symbol, and the
following empty stored block's alignment and length fields. A lower
bound skips header construction when even a minimum header cannot
improve the candidate. Only winning regions replay the deterministic
parser; no per-byte token buffer is retained. Tree nodes, sorting
buffers, traversal stacks, tables, and header runs occupy reusable
device workspace.

For \code{inputBytes} filtered bytes in \code{regionCount} regions,
stored-block fallback bounds the complete zlib output by
\code{inputBytes + 5 * regionCount + 11} bytes. A CUB device prefix
scan computes output offsets, and a compaction kernel assembles the
regions in their original order.

\subsubsection{Checksums and PNG framing}

The image uses one zlib header, one final empty stored block, and one
Adler-32 trailer \cite{RFC1950}. IDAT boundaries do not start new
compression streams. Each IDAT payload is at most 64~KiB. The GPU
generates PNG chunk headers and CRCs, combining ordered warp segments
through scalar arithmetic over the reflected CRC polynomial. The host
receives small result metadata and the final PNG byte vector.

Compression selection is automatic and requires no public mode
argument. Stronger match searching and dictionaries spanning region
boundaries remain possible future improvements. The current independent
regions favor parallelism and bounded workspace at the cost of some
compression ratio.

\subsubsection{Measured encoding cost and output size}
\label{subsec:png-benchmarks}

\Cref{tab:png-current-benchmark,tab:png-compression-sizes} report
3840~by~2160 opaque images measured on October~4, 2026 with an NVIDIA
GeForce RTX~5090 (\code{sm\_120}), CUDA Toolkit/runtime~13.2, and
Visual Studio~2026 in Windows Release. The measurements were collected
after builds and correctness tests completed, with no concurrent build
or test workload. The benchmark procedure is given in
\cref{subsec:png-validation}.

GPU timing includes serialization, filtering, compression, checksums,
PNG assembly, stream waits, and transfer of completed bytes. Input
upload and construction of the CPU \code{WPngImage} occur outside the
measured interval. Warm GPU and CPU times are medians of three calls;
first-call timing separately includes initial workspace allocation.
Fractal rendering, palette conversion, iteration upload, and disk
writing are excluded. The CPU baseline uses
\code{PngEncodingOptions\{false, false, 32\}}. These measurements
characterize the fixtures and hardware rather than end-to-end save
latency or a timing guarantee.

\begin{table}[!htbp]
\centering
\small
\setlength{\tabcolsep}{4pt}
\caption[4K PNG encoding measurements]{4K PNG encoding measurements
after production save integration. Times are warm medians in
milliseconds; workspace is retained encoder allocation capacity in MiB.}
\label{tab:png-current-benchmark}
\begin{tabularx}{\textwidth}{Xrrrrr}
\hline
Image & GPU ms & CPU ms & GPU bytes & CPU bytes & MiB \\
\hline
Uniform      & 2.42  & 502.36  & 76,027     & 51,113     & 264.3 \\
Gradient     & 5.31  & 671.04  & 3,477,715  & 1,147,577  & 267.5 \\
Seeded noise & 42.79 & 1883.63 & 49,785,331 & 49,794,895 & 311.7 \\
Mandelbrot   & 13.12 & 598.19  & 857,686    & 692,127    & 265.0 \\
\hline
\end{tabularx}
\end{table}

The Mandelbrot fixture encodes in approximately 13.12~ms on the GPU
versus 598.19~ms on the CPU, while its GPU file is approximately
1.24~times the CPU size. The gradient fixture retains a larger size
penalty because the greedy parser and independent dictionaries limit
match discovery. Noise selects stored blocks and has nearly the same
output size as the CPU encoder. Retained workspace excludes the
63.3~MiB input image and all host allocations.

\begin{table}[!htbp]
\centering
\small
\caption[PNG compression size comparison]{PNG compression size comparison
on identical filtered 4K inputs. Dynamic coding replaces the fixed/stored
candidate only when smaller; equal sizes retain that candidate.}
\label{tab:png-compression-sizes}
\begin{tabularx}{\textwidth}{Xrrr}
\hline
Image & Fixed/stored bytes & Dynamic bytes & Reduction \\
\hline
Uniform      & 325,145    & 76,027     & 76.6\% \\
Gradient     & 5,082,472  & 3,477,715  & 31.6\% \\
Seeded noise & 49,785,331 & 49,785,331 & 0\% \\
Mandelbrot   & 1,180,922  & 857,686    & 27.4\% \\
\hline
\end{tabularx}
\end{table}

Dynamic code construction and parser replay trade encoding work for
smaller files and add approximately 45~MiB to retained 4K workspace
relative to the initial fixed/stored implementation. Timing is
diagnostic; exact decoded pixels and compression-size bounds are the
benchmark's assertions.
\subsection{Imagina-Compatible Orbit Format}
\label{subsec:imagina-format}

\FractalShark{} can import and export reference orbits in a binary format
compatible with the Imagina fractal viewer, enabling cross-tool validation of
orbit computations.  The file layout is header-driven: a 32-byte file header
contains a magic number (identifying either Imagina or \FractalShark{} origin), a
reserved field, and byte offsets to the location and reference-orbit sections.
The location section stores the reference-point coordinates and zoom metadata at
arbitrary precision; the reference-orbit section stores orbit data together with
precision radii (absolute and relative) and a validity radius.

For orbits that include linear-approximation (LA) data, additional metadata
accompanies the orbit payload: the reference point in single precision,
iteration counts, periodicity status, and a flag indicating whether Approximation
Transform (AT) acceleration is enabled.  AT parameters comprise the step
length, escape radius, and LA coefficients (\code{ZCoeff}, \code{CCoeff},
\code{InvZCoeff}), as well as the number of LA stages, allowing the importing
application to reconstruct the full approximation hierarchy described in
\cref{sec:la-stage-hierarchy}.  The format therefore supports both basic
reference orbits (perturbation data only) and full LA-equipped orbits.

\FractalShark{} exposes four save modes: uncompressed text, simple compressed,
max compressed, and Imagina-compatible max compressed.  The Imagina-compatible
variant writes the orbit in Imagina's native binary layout so that the file can
be imported directly.  Loading supports two modes: \emph{Match}, which validates
that the loaded orbit's coordinates and precision match the current view, and
\emph{Use Saved}, which loads unconditionally and adopts the file's parameters.
\subsection{Cross-Platform Considerations}
\label{sec:cross-platform}

\FractalShark{} supports Windows and Linux as peer platforms. Windows builds use
Visual~Studio~2026 and MSBuild; Linux builds use CMake and Clang, with release
binaries targeted at Ubuntu. Version~0.53 introduced the native Linux target.
The Linux build includes the CPU numeric core, CUDA rendering kernels, GPU
high-precision arithmetic, the command-line renderer, the portable unit tests,
and the \code{FractalSharkGuiLinux} graphical application. The native Linux and
Win32 graphical applications are maintained at feature parity over the same
renderer and portable command contracts.

\subsubsection{Dual build system}
\label{subsec:dual-build}

The Windows working tree and Visual~Studio solution are authoritative. Shared
projects maintain parallel \code{.vcxproj} and \code{CMakeLists.txt} source
lists and settings. Projects that remain Windows-only do not have CMake targets;
the Win32 GUI and native Linux GUI are intentionally separate applications over
the same renderer and portable command contracts. Adding or removing a source
file from a shared project therefore requires updating both build descriptions.

The Linux target set comprises \code{FractalSharkPlatform},
\code{FractalSharkLib}, \code{FractalSharkGpuLib},
\code{HpSharkFloatLib}, the explicit-instantiation generator, the GPU test
harness and its support library, \code{FractalSharkTest},
\code{FractalSharkCli}, and \code{FractalSharkGuiLinux}. The GUI additionally
requires Xlib, GLX, and OpenGL. Dear~ImGui is fetched at a pinned version during
CMake configuration. Splash images and Imagina test fixtures are embedded into
ELF objects during the build, so the packaged executables do not depend on
loose copies of those assets. The splash-image embedding step currently limits
the GUI target to x86-64 ELF systems.

Continuous integration configures and builds both Debug and Release on Linux,
runs \code{FractalSharkTest}, and smoke-tests a release CLI render. It verifies
that the CLI and GUI binaries have no unresolved shared-library dependencies
and packages both release executables in the Linux archive. The GUI is built
but is not launched on the headless runner. Separate workflow jobs build and
test Windows and compile this document; release publication depends on all
three jobs.

\subsubsection{Building on Linux}
\label{subsec:building-linux}

For normal local Linux development, run \code{./build\_linux.sh}. The script
configures and builds the persistent \code{build-debug/} and
\code{build-release/} directories, which should remain in place for repeated
manual testing and incremental rebuilds. After configuration, a Debug
incremental build may use \code{cmake -{}-build build-debug -{}-parallel};
the release equivalent uses \code{build-release}. Debug validation runs
\code{./build-debug/FractalSharkTest/FractalSharkTest}.

\code{HpSharkFloatTest} requires physical CUDA hardware for its runtime
checks, so it is not expected to run on ordinary hosted CI. The GitHub Actions
workflow at \code{.github/workflows/build.yml} remains the authoritative Linux
environment setup reference. It records the exact package versions, CUDA and
compiler selection, CMake flags, LFS materialization checks, dependency audits,
release packaging behavior, and other details that should not be duplicated in
prose.

\subsubsection{Platform abstraction layer}
\label{subsec:platform-abstraction}

The \code{FractalSharkPlatform} library concentrates operating-system behavior
behind portable interfaces. The \code{Environment} namespace covers aligned
and system-heap allocation, file and directory operations, process and memory
information, timing, thread naming and affinity, console access, input polling,
temporary paths, embedded resources, diagnostics, and small user-interface
helpers. Portable geometry and opaque-handle types prevent Win32 types from
leaking into shared library interfaces.

Dedicated abstractions cover crash handling, process memory limits, wait-cursor
state, local IPC transport, and native OpenGL contexts. Each abstraction has
separate Win32 and Linux translation units. The CLI server uses the local IPC
interface with named pipes on Windows and Unix-domain sockets on Linux. The
OpenGL wrapper uses WGL on Windows and GLX on Linux while presenting the same
client-rectangle, context, buffer-swap, renderer-detection, and texture-limit
interface to the rendering code. On Windows, the crash handler writes a
\code{.dmp} minidump for post-mortem analysis.

This separation avoids distributing platform conditionals throughout the
numeric and rendering libraries. Platform-specific graphical behavior remains
in \code{FractalShark::Win32} or \code{FractalShark::Linux}; portable commands
and menu contracts remain directly in \code{FractalShark}. Operating-system
entry points and third-party callback interfaces remain in the global namespace
when required by their APIs.

\subsubsection{Native Linux graphical application}
\label{subsec:linux-gui}

\code{FractalSharkGuiLinux} uses Xlib for its top-level window and event loop,
GLX for the native OpenGL context, and the shared OpenGL display pipeline for
presenting completed frames. The GUI thread owns the window, GLX context, and
all Dear~ImGui state. It forwards input to an in-project ImGui Xlib backend,
presents the latest renderer output, draws overlays, and swaps buffers from the
same event loop. The backend handles mouse buttons and wheels, keyboard input,
focus changes, window resizing, and clipboard integration.

The right-click menu is a native X11 popup hierarchy rather than an ImGui
window. It is generated from the shared \code{MenuTree} catalog and dispatches
through the same portable command identifiers as the Win32 GUI. A shared
\code{MenuState} implementation supplies enablement, check marks, and radio
selections to both native menu builders. Linux-specific command handlers
translate commands that require X11 operations. This arrangement keeps algorithm selection,
navigation, built-in views, iteration controls, perturbation modes, palettes,
tests, benchmarks, and linear-approximation settings synchronized with the
portable command catalog while leaving window-system operations in the Linux
front end. Palette rotation is handled by the portable command layer and uses
the shared render-pool presentation path to recolor the current iteration
buffer.

Dear~ImGui supplies the overlay surfaces that do not need separate native
windows: help and rendering-information messages, file selection, saved-view
selection, manual location entry, and the drag-zoom rectangle. The X11 front
end also implements resize-driven recalculation, mouse-wheel zoom, captured
drag zoom, fullscreen and square-fullscreen transitions, minimization, feature
finding and automatic zoom, save and load commands, and an X11 clipboard owner
for UTF-8 text. A separate Xlib splash window displays a randomly selected
embedded image while initialization proceeds.

The Linux GUI is functional, and its feature set should continue to converge on
the Win32 GUI. Current limitations should be stated explicitly as implementation
gaps rather than hidden behind a premature claim of completed parity.

\subsubsection{Arbitrary-precision serialization}
\label{subsec:mpir-portability}

\FractalShark{} relies on MPIR for arbitrary-precision integer arithmetic on
Windows and on GMP for the same role on Linux. The two libraries are
source-compatible at the API level but differ in some internal representation
details, which matters because reference-orbit files and saved views contain
serialized arbitrary-precision integers that must round-trip across platforms.

To guarantee byte-identical files, \FractalShark{} defines a single canonical
wire format that matches MPIR's native \code{\_\_mpz\_struct} layout. On
Windows this format is produced by direct field access into the underlying
MPIR object. On Linux the same byte sequence is reconstructed atop GMP using
\code{mpz\_export} and \code{mpz\_import}, which expose limb data through a
portable interface. The two implementations are validated against a checked-in
golden file (\code{golden\_mpir.bin}) by the \code{TestMpirSerialization}
unit test, which runs on both platforms in CI. Any divergence in the wire
format fails the test on whichever platform regressed, preventing silent
incompatibility between Windows-written and Linux-written orbit files.
